\documentclass[tikz,border=6pt]{standalone}
% Shared visual language for the thesis architecture figures.
% Compile each standalone source with XeLaTeX from this directory.
\usepackage[UTF8,scheme=plain,fontset=none]{ctex}
\usepackage{amsmath,amssymb}
\usepackage{fontspec}
\setmainfont{TeX Gyre Termes}
\setsansfont{TeX Gyre Heros}
\setmonofont{DejaVu Sans Mono}
\setCJKmainfont{Noto Serif SC}
\setCJKsansfont[BoldFont=SimHei]{Noto Sans SC}
\usetikzlibrary{arrows.meta,positioning,calc,fit,backgrounds,shapes.geometric,matrix,decorations.pathreplacing}

\definecolor{Ink}{HTML}{203040}
\definecolor{Muted}{HTML}{66788A}
\definecolor{LineSoft}{HTML}{AAB7C4}
\definecolor{DataBlue}{HTML}{2563A6}
\definecolor{DataFill}{HTML}{EAF3FB}
\definecolor{ClockPurple}{HTML}{7450A5}
\definecolor{ClockFill}{HTML}{F1ECF8}
\definecolor{ControlTeal}{HTML}{16827B}
\definecolor{ControlFill}{HTML}{E9F6F4}
\definecolor{DspPurple}{HTML}{6A4C93}
\definecolor{DspFill}{HTML}{F0EAF7}
\definecolor{RamGreen}{HTML}{367C59}
\definecolor{RamFill}{HTML}{ECF6F0}
\definecolor{FabricOrange}{HTML}{C56A1A}
\definecolor{FabricFill}{HTML}{FFF1E5}
\definecolor{EvidenceGreen}{HTML}{2F7D4B}
\definecolor{EvidenceFill}{HTML}{EAF6EE}
\definecolor{Amber}{HTML}{B07A16}
\definecolor{AmberFill}{HTML}{FFF6DD}
\definecolor{PendingGray}{HTML}{7C8792}
\definecolor{PendingFill}{HTML}{F1F3F5}
\definecolor{Danger}{HTML}{A04444}
\definecolor{DangerFill}{HTML}{FFF0F0}

\tikzset{
  >={Stealth[length=2.2mm,width=1.45mm]},
  every node/.style={font=\sffamily\footnotesize,text=Ink},
  data/.style={-{Stealth[length=2.2mm,width=1.45mm]},draw=DataBlue,line width=0.92pt},
  data bi/.style={{Stealth[length=1.9mm,width=1.3mm]}-{Stealth[length=1.9mm,width=1.3mm]},draw=DataBlue,line width=0.82pt},
  control/.style={-{Stealth[length=2.0mm,width=1.35mm]},draw=ControlTeal,dashed,line width=0.78pt},
  clock/.style={-{Stealth[length=2.0mm,width=1.35mm]},draw=ClockPurple,densely dotted,line width=0.9pt},
  address/.style={-{Stealth[length=2.0mm,width=1.35mm]},draw=FabricOrange,dash dot,line width=0.76pt},
  feedback/.style={-{Stealth[length=1.8mm,width=1.25mm]},draw=Muted,line width=0.66pt},
  block/.style={draw=Ink,fill=white,rounded corners=1.5pt,line width=0.68pt,minimum height=10mm,minimum width=23mm,align=center,inner sep=3pt},
  data block/.style={block,draw=DataBlue,fill=DataFill},
  control block/.style={block,draw=ControlTeal,fill=ControlFill},
  clock block/.style={block,draw=ClockPurple,fill=ClockFill},
  dsp/.style={block,draw=DspPurple,fill=DspFill,line width=0.9pt},
  ram/.style={block,draw=RamGreen,fill=RamFill,line width=0.9pt},
  fabric/.style={block,draw=FabricOrange,fill=FabricFill,line width=0.9pt},
  result/.style={block,draw=EvidenceGreen,fill=EvidenceFill,line width=0.9pt},
  pending/.style={block,draw=PendingGray,fill=PendingFill,line width=0.8pt},
  warning/.style={block,draw=Amber,fill=AmberFill,line width=0.9pt},
  boundary/.style={draw=Muted!75,rounded corners=2.2mm,dashed,line width=0.55pt,inner sep=3mm},
  title/.style={font=\sffamily\bfseries\large,text=Ink,align=center},
  subtitle/.style={font=\sffamily\scriptsize,text=Muted,align=center},
  group title/.style={font=\sffamily\bfseries\footnotesize,text=Ink,anchor=west},
  signal/.style={font=\sffamily\scriptsize,text=Ink,fill=white,inner sep=1.2pt,align=center},
  note/.style={font=\sffamily\scriptsize,text=Muted,align=center},
  tiny note/.style={font=\sffamily\tiny,text=Muted,align=center},
  badge/.style={draw=EvidenceGreen,fill=EvidenceFill,rounded corners=1.8mm,line width=0.65pt,font=\sffamily\bfseries\scriptsize,text=EvidenceGreen,inner xsep=4pt,inner ysep=2pt,align=center},
  pending badge/.style={draw=PendingGray,fill=PendingFill,rounded corners=1.8mm,line width=0.65pt,font=\sffamily\bfseries\scriptsize,text=PendingGray,inner xsep=4pt,inner ysep=2pt,align=center},
  innovation/.style={circle,draw=Amber,fill=AmberFill,line width=0.8pt,minimum size=6.5mm,font=\sffamily\bfseries\scriptsize,text=Amber,inner sep=0pt},
  sum/.style={circle,draw=Ink,fill=white,line width=0.75pt,minimum size=6mm,inner sep=0pt},
  mux/.style={trapezium,trapezium left angle=72,trapezium right angle=108,draw=Ink,fill=white,line width=0.68pt,minimum height=11mm,minimum width=12mm,align=center,inner sep=2pt}
}

\newcommand{\bitrate}[2]{{\scriptsize #1 bit，#2}}
\newcommand{\passmark}{\textcolor{EvidenceGreen}{\bfseries 满足判据}}
\newcommand{\pendingmark}{\textcolor{PendingGray}{\bfseries 待验证}}


\begin{document}
\begin{tikzpicture}[x=1cm,y=1cm]

\node[title] at (12.5,8.45) {24/20/20 bit三级FIR--CIC插值链及有效事件率演化};
\node[subtitle] at (12.5,8.00)
  {当前2-DSP目标实现配置；补偿并入第三级FIR系数；全链采用单一$128F_s$音频主时钟};

% Main data path.
\node[data block,minimum width=21mm,minimum height=15mm] (xin) at (0.95,5.20)
  {$x[n]$\\[-1pt]\bitrate{24}{$F_s$}};
\node[dsp,minimum width=35mm,minimum height=25mm] (s1) at (4.00,5.20)
  {第一级半带FIR $\times2$\\[-1pt]
   {\scriptsize 105抽头（104阶）}\\[-1pt]
   {\scriptsize 52历史字顺序MAC}\\[-1pt]
   {\scriptsize ACC41 / Q15}\\[-1pt]
   {\scriptsize DSP48E1\#0；$y_2$:24 bit@$2F_s$}};
\node[fabric,minimum width=26mm,minimum height=20mm] (b12) at (7.25,5.20)
  {量化/相位桥接\\[-1pt]
   {\scriptsize 24$\rightarrow$20 bit}\\[-1pt]
   {\scriptsize $\gg4$舍入饱和}\\[-1pt]
   {\scriptsize pending/phase}};
\node[data block,minimum width=31mm,minimum height=22mm] (s2) at (10.55,5.20)
  {第二级多相FIR $\times2$\\[-1pt]
   {\scriptsize $P_0/P_1$: 9/8 MAC}\\[-1pt]
   {\scriptsize 20$\rightarrow$20 bit，ACC38}\\[-1pt]
   {\scriptsize $y_4$@$4F_s$；板级$4\times$抽头}};
\node[fabric,minimum width=23mm,minimum height=18mm] (b23) at (13.45,5.20)
  {相位桥接\\[-1pt]{\scriptsize 20$\rightarrow$20 bit}\\[-1pt]{\scriptsize 无再量化}};
\node[data block,minimum width=36mm,minimum height=27mm] (s3) at (16.65,5.20)
  {第三级多相FIR $\times2$\\[-1pt]
   {\scriptsize $P_0/P_1$: 6/5 MAC}\\[-1pt]
   {\scriptsize flat / P3补偿系数库}\\[-1pt]
   {\scriptsize $y_{8,int}$:21 bit@$8F_s$}\\[-1pt]
   {\scriptsize 板级$8\times$: sat20$\rightarrow\ll4$}};
\node[fabric,minimum width=34mm,minimum height=27mm] (cic) at (20.70,5.20)
  {等价CIC $\times16$\\[-1pt]
   {\scriptsize $C^2\rightarrow\mathrm{Hold}_{16}\rightarrow I^2$}\\[-1pt]
   {\scriptsize 21$\rightarrow$23$\rightarrow$26$\rightarrow$29 bit}\\[-1pt]
   {\scriptsize $\gg8$舍入饱和；0 DSP}\\[-1pt]
   {\scriptsize $y_{128}$:20 bit@$128F_s$}};
\node[data block,minimum width=24mm,minimum height=18mm] (restore) at (24.00,5.20)
  {$Q$格式恢复\\[-1pt]{\scriptsize 20$\rightarrow$24 bit}\\[-1pt]{\scriptsize 左移4位}};

\draw[data] (xin) -- (s1);
\draw[data] (s1) -- (b12);
\draw[data] (b12) -- (s2);
\draw[data] (s2) -- (b23);
\draw[data] (b23) -- (s3);
\draw[data] (s3) -- (cic);
\draw[data] (cic) -- (restore);

% Shared Stage2/3 physical engine.
\node[dsp,minimum width=54mm,minimum height=13mm] (sharedDsp) at (13.65,2.65)
  {DSP48E1\#1：Stage2/Stage3时分共享MAC、Q15舍入与饱和};
\draw[address] (s2.south) -- node[signal,left] {job $S_2$} (sharedDsp.north west);
\draw[address] (s3.south) -- node[signal,right] {job $S_3$} (sharedDsp.north east);
\node[ram,minimum width=37mm,minimum height=13mm] (hist1) at (4.00,2.65)
  {RAMB18E1\#1\\[-1pt]{\scriptsize Stage1历史}};
\node[ram,minimum width=37mm,minimum height=13mm] (hist23) at (9.30,1.15)
  {RAMB18E1\#2\\[-1pt]{\scriptsize Stage2/3统一历史}};
\node[ram,minimum width=43mm,minimum height=13mm] (coeff) at (17.15,1.15)
  {RAMB18E1\#3\\[-1pt]{\scriptsize 三级FIR统一系数}};
\draw[address] (hist1) -- (s1.south);
\draw[address] (hist23.north) -- (sharedDsp.south west);
\draw[address] (coeff.north) -- (sharedDsp.south east);

% CE and valid plane.
\node[control block,minimum width=30mm] (ce) at (1.55,7.15) {CE/phase发生器\\[-1pt]{\scriptsize 单一$128F_s$时钟}};
\draw[control] (ce.east) -- ++(1.0,0) -| (s1.north);
\draw[control] (ce.east) -- ++(2.0,0) -| (s2.north);
\draw[control] (ce.east) -- ++(3.0,0) -| (s3.north);
\draw[control] (ce.east) -- ++(4.0,0) -| (cic.north);

% Output taps and innovation statements.
\draw[feedback] (xin.south) -- ++(0,-0.92) node[note,below] {$1\times$原始PCM抽头};
\draw[feedback] (restore.south) -- ++(0,-0.92) node[note,below] {$128\times$可见抽头};

\node[innovation] at (18.75,7.62) {I1};
\node[note,anchor=west,align=left] at (19.15,7.62) {CIC通带补偿折叠至\\Stage3系数库};
\node[innovation] at (22.20,7.62) {I2};
\node[note,anchor=west,align=left] at (22.60,7.62) {保持型CIC等价变换\\消除一级积分/梳状状态};

\begin{scope}[on background layer]
  \node[boundary,fill=DataFill!35,fit=(s2)(b23)(s3)(sharedDsp)(hist23)(coeff)] (sharedGroup) {};
\end{scope}
\node[group title] at ($(sharedGroup.north west)+(0.20,-0.25)$) {Stage2/Stage3共享FIR引擎};

\node[note,anchor=west] at (0.35,0.15)
  {核心资源：2 DSP48E1 + 3 RAMB18E1；CIC的差分器和两级积分器均映射至Fabric/CARRY。};
\node[note,anchor=east] at (24.85,0.15)
  {粗实线：样值数据；青色虚线：CE/phase；橙色点划线：任务与存储器访问。};

\end{tikzpicture}
\end{document}


